\section{証明の範囲と再現}

\subsection{確定した一般現象と残る分類}

全整数 $m\ge3$ で、引戻し成分の一般点のArtin環
$K[a,b]/(a^{m+2},ab,b^2)$ と交叉環を決定した。
接線系列では、$r=m-2j>0$ の全範囲で既約成分性と
全形式局所環 $(\kappa u,\gamma^{j+1}u)$ を証明した。
一般点の長さ $j+1$ は成分に固有であり、固定した重みでも互いに異なる長さが現れる。
$m=1,2$ には二重ベクトル場の追加係数があるので、同じ成分定理の対象には含めない。
接線成分定理の $r>0$ と、量子化・中心の定理の $r\ge0$ も区別する。

全重みの非線形アフィン族について、Nakayama自己同型、通常のCY条件、
次数付き自己同型群、多項式中心、次数零の関数斜体を決定した。
多項式中心がDu Val曲面であるという結論は、その曲面の座標環を中心として含むという意味であり、
曲面自身に非可換な積を入れたという意味ではない。
非零の量子化係数 $s$ に対するこれらの代数は非PIであり、中心上有限な代数として扱わない。
同じ重みの異なる共鳴の射影圏が同値かどうかは未決である。

重み三では少なくとも十成分を認定したが、全分類と最大次元は未決である。
$c=0$ の被約部分の三成分とその大域所属、非零の接線型線形場に対する
四共鳴の形式近傍は決定した。
残る層には非被約三次形式と、その他の特異三次形式の特殊係数がある。
混合核の増大だけから新成分を数えることはしない。

三角形の三つの直線・二次曲線枝は、同じ $L(1,2,3)$ 成分に属する。
共点三直線の交叉定理は、明記した非空開集合に対するものであり、
全境界点へ同じ局所式を延長する主張ではない。
全ての次元は自己同型群とスカラー倍で商する前の値である。

\subsection{構造的証明と有限証明書}

本版の統合検証器は \nolinkurl{code/verify_v34.py} である。
全重みの共鳴に関する記号恒等式、重み三の全Schur式、
Nakayamaの規約を確認する三次syzygy、不変環の関係式・等変変換と中心式を検証する。
通常実行では入力や証明書を再生成せず、全JSON出力を保存済み証明書と比較する。
コード・入力・証明書の実行前後のhashも比較する。
新しい定理の独立監査を保存し、過去の原稿・計算入力・検証記録を保持する。

全重みの結論は有限個の検算からの外挿ではない。
混合核の二重対角表示、横断二次曲線の次数降下、余核指標、
特異点評価が、任意の $m,j$ でSchurイデアルを決める。
不変環の生成性は局所化と $x$ 飽和で、中心の生成性は全ての零重量単項式の分類で証明する。
AS正則性、CY性、例外列の存在は、本文の構造的定理と適用条件による。
補助的な有限次数の核次元検査は、これらの生成性の証明を置き換えない。

miniversal基底の同定は形式座標と量子化の選択を伴う。
固定した非零量子代数の全変形空間、全Poisson成分のAS正則量子化、
元の四対象cornerの全変形についての量子座標環の実現可能性は、本稿では決定していない。
$m>3$ の四対象cornerが無障害であることや、外側との $HH^2$ 制限が同型であることも未証明である。
重み三の標準補完に関する\cite{VKAmbient32}の主張を、
別の例外列や一般重みに拡張しない。
